%%%%%%%%%%%%%%%%%%%%%%%%%%%%%%%%%%%%%%%%%%%%%%%%%%%%%%%%%%%%%%%%%%%%%%%%%%%%%%%%
%%%%%%%%%%%%%%%%%%%%%%%%%%%%%%%%%%%%%%%%%%%%%%%%%%%%%%%%%%%%%%%%%%%%%%%%%%%%%%%%
%%                                                                            %%
%% thesistemplate.tex version 3.20 (2018/08/31)                               %%
%% The LaTeX template file to be used with the aaltothesis.sty (version 3.20) %%
%% style file.                                                                %%
%% This package requires pdfx.sty v. 1.5.84 (2017/05/18) or newer.            %%
%%                                                                            %%
%% This is licensed under the terms of the MIT license below.                 %%
%%                                                                            %%
%% Written by Luis R.J. Costa.                                                %%
%% Currently developed at the Learning Services of Aalto University School of %%
%% Electrical Engineering by Luis R.J. Costa since May 2017.                  %%
%%                                                                            %%
%% Copyright 2017-2018, by Luis R.J. Costa, luis.costa@aalto.fi,              %%
%% Copyright 2017-2018 Swedish translations in aaltothesis.cls by Elisabeth   %%
%% Nyberg, elisabeth.nyberg@aalto.fi and Henrik Wallén,                       %%
%% henrik.wallen@aalto.fi.                                                    %%
%% Copyright 2017-2018 Finnish documentation in the template opinnatepohja.tex%%
%% by Perttu Puska, perttu.puska@aalto.fi, and Luis R.J. Costa.               %%
%% Copyright 2018 English template thesistemplate.tex by Luis R.J. Costa.     %%
%% Copyright 2018 Swedish template kandidatarbetsbotten.tex by Henrik Wallen. %%
%%                                                                            %%
%% Permission is hereby granted, free of charge, to any person obtaining a    %%
%% copy of this software and associated documentation files (the "Software"), %%
%% to deal in the Software without restriction, including without limitation  %%
%% the rights to use, copy, modify, merge, publish, distribute, sublicense,   %%
%% and/or sell copies of the Software, and to permit persons to whom the      %%
%% Software is furnished to do so, subject to the following conditions:       %%
%% The above copyright notice and this permission notice shall be included in %%
%% all copies or substantial portions of the Software.                        %%
%% THE SOFTWARE IS PROVIDED "AS IS", WITHOUT WARRANTY OF ANY KIND, EXPRESS OR %%
%% IMPLIED, INCLUDING BUT NOT LIMITED TO THE WARRANTIES OF MERCHANTABILITY,   %%
%% FITNESS FOR A PARTICULAR PURPOSE AND NONINFRINGEMENT. IN NO EVENT SHALL    %%
%% THE AUTHORS OR COPYRIGHT HOLDERS BE LIABLE FOR ANY CLAIM, DAMAGES OR OTHER %%
%% LIABILITY, WHETHER IN AN ACTION OF CONTRACT, TORT OR OTHERWISE, ARISING    %%
%% FROM, OUT OF OR IN CONNECTION WITH THE SOFTWARE OR THE USE OR OTHER        %%
%% DEALINGS IN THE SOFTWARE.                                                  %%
%%                                                                            %%
%%                                                                            %%
%%%%%%%%%%%%%%%%%%%%%%%%%%%%%%%%%%%%%%%%%%%%%%%%%%%%%%%%%%%%%%%%%%%%%%%%%%%%%%%%
%%                                                                            %%
%%                                                                            %%
%% An example for writting your thesis using LaTeX                            %%
%% Original version and development work by Luis Costa, changes to the text   %% 
%% in the Finnish template by Perttu Puska.                                   %%
%% Support for Swedish added 15092014                                         %%
%% PDF/A-b support added on 15092017                                          %%
%% PDF/A-2 support added on 24042018                                          %%
%%                                                                            %%
%% This example consists of the files                                         %%
%%         thesistemplate.tex (version 3.20) (for text in English)            %%
%%         opinnaytepohja.tex (version 3.20) (for text in Finnish)            %%
%%         kandidatarbetsbotten.tex (version 1.00) (for text in Swedish)      %%
%%         aaltothesis.cls (versio 3.20)                                      %%
%%         kuva1.eps (graphics file)                                          %%
%%         kuva2.eps (graphics file)                                          %%
%%         kuva1.jpg (graphics file)                                          %%
%%         kuva2.jpg (graphics file)                                          %%
%%         kuva1.png (graphics file)                                          %%
%%         kuva2.png (graphics file)                                          %%
%%         kuva1.pdf (graphics file)                                          %%
%%         kuva2.pdf (graphics file)                                          %%
%%                                                                            %%
%%                                                                            %%
%% Typeset in Linux either with                                               %%
%% pdflatex: (recommended method)                                             %%
%%             $ pdflatex thesistemplate                                      %%
%%             $ pdflatex thesistemplate                                      %%
%%                                                                            %%
%%   The result is the file thesistemplate.pdf that is PDF/A compliant, if    %%
%%   you have chosen the proper \documenclass options (see comments below)    %%
%%   and your included graphics files have no problems.
%%                                                                            %%
%% Or                                                                         %%
%% latex: (this method is not recommended)                                    %%
%%             $ latex thesistemplate                                         %%
%%             $ latex thesistemplate                                         %%
%%                                                                            %%
%%   The result is the file thesistemplate.dvi, which is converted to ps      %%
%%   format as follows:                                                       %%
%%                                                                            %%
%%             $ dvips thesistemplate -o                                      %%
%%                                                                            %%
%%   and then to pdf as follows:                                              %%
%%                                                                            %%
%%             $ ps2pdf thesistemplate.ps                                     %%
%%                                                                            %%
%%   This pdf file is not PDF/A compliant. You must must make it so using,    %%
%%   e.g., Acrobat Pro or PDF-XChange.                                        %%
%%                                                                            %%
%%                                                                            %%
%% Explanatory comments in this example begin with the characters %%, and     %%
%% changes that the user can make with the character %                        %%
%%                                                                            %%
%%%%%%%%%%%%%%%%%%%%%%%%%%%%%%%%%%%%%%%%%%%%%%%%%%%%%%%%%%%%%%%%%%%%%%%%%%%%%%%%
%%%%%%%%%%%%%%%%%%%%%%%%%%%%%%%%%%%%%%%%%%%%%%%%%%%%%%%%%%%%%%%%%%%%%%%%%%%%%%%%
%%
%% WHAT is PDF/A
%%
%% PDF/A is the ISO-standardized version of the pdf. The standard's goal is to
%% ensure that he file is reproducable even after a long time. PDF/A differs
%% from pdf in that it allows only those pdf features that support long-term
%% archiving of a file. For example, PDF/A requires that all used fonts are
%% embedded in the file, whereas a normal pdf can contain only a link to the
%% fonts in the system of the reader of the file. PDF/A also requires, among
%% other things, data on colour definition and the encryption used.
%% Currently three PDF/A standards exist:
%% PDF/A-1: based on PDF 1.4, standard ISO19005-1, published in 2005.
%%          Includes all the requirements essential for long-term archiving.
%% PDF/A-2: based on PDF 1.7, standard ISO19005-2, published in 2011.
%%          In addition to the above, it supports embedding of OpenType fonts,
%%          transparency in the colour definition and digital signatures.
%% PDF/A-3: based on PDF 1.7, standard ISO19005-3, published in 2012.
%%          Differs from the above only in that it allows embedding of files in
%%          any format (e.g., xml, csv, cad, spreadsheet or wordprocessing
%%          formats) into the pdf file.
%% PDF/A-1 files are not necessarily PDF/A-2 -compatible and PDF/A-2 are not
%% necessarily PDF/A-1 -compatible.
%% All of the above PDF/A standards have two levels:
%% b: (basic) requires that the visual appearance of the document is reliably
%%    reproduceable.
%% a (accessible) in addition to the b-level requirements, specifies how
%%   accessible the pdf file is to assistive software, say, for the physically
%%   impaired.
%% For more details on PDF/A, see, e.g., https://en.wikipedia.org/wiki/PDF/A
%%
%%
%% WHICH PDF/A standard should my thesis conform to?
%%
%% Primarily to the PDF/A-1b standard. All the figures and graphs typically
%% use in thesis work do not require transparency features, a basic '2-D'
%% visualisation suffices. The font to be used are specified in this template
%% and they should not be changed. However, if you have figures where
%% transparency characteristics matter, use the PDF/A-2b standard. Do not use
%% the PDF/A-3b standard for your thesis.
%%
%%
%% WHAT graphics format can I use to produce my PDF/A compliant file?
%%
%% When using pdflatex to compile your work, use jpg, png or pdf files. You may
%% have PDF/A compliance problems with figures in pdf format. Do not use PDF/A
%% compliant graphics files.
%% If you decide to use latex to compile your work, the only acceptable file
%% format for your figure is eps. DO NOT use the ps format for your figures.

%% USE one of these:
%% * the first when using pdflatex, which directly typesets your document in the
%%   chosen pdf/a format and you want to publish your thesis online,

%% * the second when you want to print your thesis to bind it, or
%% * the third when producing a ps file and a pdf/a from it.
%%
\documentclass[UKenglish, 12pt, a4paper, sci, utf8, a-1b, online]{aaltothesis}
%\documentclass[english, 12pt, a4paper, elec, utf8, a-1b]{aaltothesis}
%\documentclass[english, 12pt, a4paper, elec, dvips, online]{aaltothesis}

%% Use the following options in the \documentclass macro above:
%% your school: arts, biz, chem, elec, eng, sci
%% the character encoding scheme used by your editor: utf8, latin1
%% thesis language: english, finnish, swedish
%% make an archiveable PDF/A-1b or PDF/A-2b compliant file: a-1b, a-2b
%%                    (with pdflatex, a normal pdf containing metadata is
%%                     produced without the a-*b option)
%% typeset in symmetric layout and blue hypertext for online publication: online
%%            (no option is the default, resulting in a wide margin on the
%%             binding side of the page and black hypertext)
%% two-sided printing: twoside (default is one-sided printing)
%%

%% Use one of these if you write in Finnish (see the Finnish template
%% opinnaytepohja.tex)
%\documentclass[finnish, 12pt, a4paper, elec, utf8, a-1b, online]{aaltothesis}
%\documentclass[finnish, 12pt, a4paper, elec, utf8, a-1b]{aaltothesis}
%\documentclass[finnish, 12pt, a4paper, elec, dvips, online]{aaltothesis}

\usepackage{graphicx}
\usepackage{mathtools}
%% Math fonts, symbols, and formatting; these are usually needed
\usepackage{amsfonts,amssymb,amsbsy,amsmath,amsthm,bm}
\usepackage{algorithm,algpseudocode}
\usepackage[labelfont=bf]{caption}
\usepackage{siunitx}
\sisetup{math-micro=\text{µ},text-micro=µ}
%% Change the school field to specify your school if the automatically set name
%% is wrong
% \university{aalto-yliopisto}
% \school{Sähkötekniikan korkeakoulu}

%% Edit to conform to your degree programme
%%
\degreeprogram{Mathematics and Operations Research}
%%

%% Your major
%%
\major{Mathematics}
%%

%% Major subject code
%%
\code{SCI3054}
%%
 
%% Choose one of the three below
%%
%\univdegree{BSc}
\univdegree{MSc}
%\univdegree{Lic}
%%

%% Your name (self explanatory...)
%%
\thesisauthor{Tuukka Himanka}
%%

%% Your thesis title comes here and possibly again together with the Finnish or
%% Swedish abstract. Do not hyphenate the title, and avoid writing too long a
%% title. Should LaTeX typeset a long title unsatisfactorily, you mght have to
%% force a linebreak using the \\ control characters.
%% In this case...
%% Remember, the title should not be hyphenated!
%% A possible "and" in the title should not be the last word in the line, it
%% begins the next line.
%% Specify the title again without the linebreak characters in the optional
%% argument in box brackets. This is done because the title is part of the 
%% metadata in the pdf/a file, and the metadata cannot contain linebreaks.
%%
\thesistitle{Physics-Informed Neural Networks in Probabilistic Spatio-Temporal Modelling}
%\thesistitle[Title of the thesis]{Title of\\ the thesis}
%%

%%
\place{Espoo}
%%

%% The date for the bachelor's thesis is the day it is presented
%%
\date{14.2.2024}
%%

%% Thesis supervisor
%% Note the "\" character in the title after the period and before the space
%% and the following character string.
%% This is because the period is not the end of a sentence after which a
%% slightly longer space follows, but what is desired is a regular interword
%% space.
%%
\supervisor{D.Sc.\ (Tech.)\ Harri Hakula}
%%

%% Advisor(s)---two at the most---of the thesis. Check with your supervisor how
%% many official advisors you can have.
%%
\advisor{Prof.\ Marko Laine}
\advisor{PhD Olle Räty}
%%

%% Aaltologo: syntax:
%% \uselogo{aaltoRed|aaltoBlue|aaltoYellow|aaltoGray|aaltoGrayScale}{?|!|''}
%% The logo language is set to be the same as the thesis language.
%%
\uselogo{aaltoRed}{''}
%%

%% The English abstract:
%% All the details (name, title, etc.) on the abstract page appear as specified
%% above.
%% Thesis keywords:
%% Note! The keywords are separated using the \spc macro
%%
\keywords{Physics-informed neural networks\spc Uncertainty quantification\spc Bayesian filtering}
%%

%% The abstract text. This text is included in the metadata of the pdf file as well
%% as the abstract page.
%%
\thesisabstract{
    Physics-informed neural networks have proved themselves as valuable supplements to classical mathematical models.
    Compared to regular physics-free neural networks, boundary conditions can be better enforced in the physics-informed variants.
    The computational efficiency of neural networks makes them feasible for many applications not possible with classical methods.
    Most real-world observations and mathematical models are subject to uncertainties.
    The Kalman filter is a solution to the uncertainty quantification to problems of uncertain observations and models.
    The ensemble Kalman filter provides an efficient Monte Carlo approximation when dealing with large data vectors.
    In this thesis, a neural network is constructed to infer velocity fields within the advection equation.
    The model is trained with meteorological cloud data.
    Despite the intrinsic difficulties of the training data, the neural network model is shown to learn the dynamics of the advection equation even with simpler synthetic data.
    In addition, the neural network presents comparable results compared to a commonly used method for optical flow estimation.
    The network is used as a state process operator in the ensemble Kalman filter.
    A simulated setting of partial observations with an attempt to produce complete state ensembles is tested.
    The network functions expectedly within the Kalman filter model. 
}

%% Copyright text. Copyright of a work is with the creator/author of the work
%% regardless of whether the copyright mark is explicitly in the work or not.
%% You may, if you wish, publish your work under a Creative Commons license (see
%% creaticecommons.org), in which case the license text must be visible in the
%% work. Write here the copyright text you want. It is written into the metadata
%% of the pdf file as well.
%% Syntax:
%% \copyrigthtext{metadata text}{text visible on the page}
%% 
%% In the macro below, the text written in the metadata must have a \noexpand
%% macro before the \copyright special character, and macros (\copyright and
%% \year here) must be separated by the \ character (space chacter) from the
%% text that follows. The macros in the argument of the \copyrighttext macro
%% automatically insert the year and the author's name. (Note! \ThesisAuthor is
%% an internal macro of the aaltothesis.cls class file).
%% Of course, the same text could have simply been written as
%% \copyrighttext{Copyright \noexpand\copyright\ 2018 Eddie Engineer}
%% {Copyright \copyright{} 2018 Eddie Engineer}
%%
\copyrighttext{Copyright \noexpand\copyright\ \number\year\ \ThesisAuthor}
{Copyright \copyright{} \number\year{} \ThesisAuthor}

%% You can prevent LaTeX from writing into the xmpdata file (it contains all the 
%% metadata to be written into the pdf file) by setting the writexmpdata switch
%% to 'false'. This allows you to write the metadata in the correct format
%% directly into the file thesistemplate.xmpdata.
%\setboolean{writexmpdatafile}{false}
\newcommand{\inputpypgf}[1]{%
  \InputIfFileExists{#1}{}{\typeout{No file #1.}}%
}
\newcommand*\mean[1]{\overline{#1}}
\newcommand{\R}{\mathbb{R}}
\newcommand{\Z}{\mathbb{Z}}
\newcommand{\N}{\mathbb{N}}
\newtheorem{theorem}{Theorem}
\newtheorem{lemma}{Lemma}
\newtheorem{definition}{Definition}
\numberwithin{equation}{section}
\numberwithin{figure}{section}
\numberwithin{theorem}{section}
\numberwithin{definition}{section}
\numberwithin{lemma}{section}

\usepackage{tikz}
\usetikzlibrary{arrows,positioning}
\def\inputArray{{1,2,5,4,5,2,4,3,2,1,6,1,0,1,5,3,1,4,1,2,5,3,1,3,5}}
\def\kernelArray{{1,1,0,0,1,1,0,1,0}}
%\usepackage{pgfplots}
%\usepgfplotslibrary{external} 
%\tikzexternalize

%% All that is printed on paper starts here
%%
\begin{document}

%% Create the coverpage
%%
\makecoverpage

%% Typeset the copyright text.
%% If you wish, you may leave out the copyright text from the human-readable
%% page of the pdf file. This may seem like a attractive idea for the printed
%% document especially if "Copyright (c) yyyy Eddie Engineer" is the only text
%% on the page. However, the recommendation is to print this copyright text.
%%
\makecopyrightpage

%% Note that when writting your thesis in English, place the English abstract
%% first followed by the possible Finnish or Swedish abstract.

%% Abstract text
%% All the details (name, title, etc.) on the abstract page appear as specified
%% above.
%%
\begin{abstractpage}[english]
    Physics-informed neural networks have proved themselves as valuable supplements to classical mathematical models.
    Compared to regular physics-free neural networks, boundary conditions can be better enforced in the physics-informed variants.
    The computational efficiency of neural networks makes them feasible for many applications not possible with classical methods.

    Most real-world observations and mathematical models are subject to uncertainties.
    The Kalman filter is a solution to the uncertainty quantification to problems of uncertain observations and models.
    The ensemble Kalman filter provides an efficient Monte Carlo approximation when dealing with large data vectors.

    In this thesis, a neural network is constructed to infer velocity fields within the advection equation.
    The model is trained with meteorological cloud data.
    Despite the intrinsic difficulties of the training data, the neural network model is shown to learn the dynamics of the advection equation even with simpler synthetic data.
    In addition, the neural network presents comparable results compared to a commonly used method for optical flow estimation.

    The network is used as a state process operator in the ensemble Kalman filter.
    A simulated setting of partial observations with an attempt to produce complete state ensembles is tested.
    The network functions expectedly within the Kalman filter model. 
\end{abstractpage}

%% The text in the \thesisabstract macro is stored in the macro \abstractext, so
%% you can use the text metadata abstract directly as follows:
%%
%\begin{abstractpage}[english]
%	\abstracttext{}
%\end{abstractpage}

%% Force a new page so that the possible Finnish or Swedish abstract does not
%% begin on the same page
%%
\newpage
%%
%% Abstract in Finnish.  Delete if you don't need it. 
%%
\thesistitle{Fysiikkainformoidut neuroverkot todennäköisyyspohjaisessa spatiotemporaalisessa mallinnuksessa}
\supervisor{TkT Harri Hakula}
\advisor{Prof. Marko Laine, FT Olle Räty}
\degreeprogram{Mathematics and Operations Researach}
\department{Matematiikan ja systeemianalyysin laitos}
\major{Mathematics}
%% The keywords need not be separated by \spc now.
\keywords{Fysiikkainformoitu neuroverkko, epävarmuuden määritys, Bayesilainen suodatus}
%% Abstract text
\begin{abstractpage}[finnish]
    Fysiikkainformoidut neuroverkot ovat osoittautuneet hyödyllisiksi työkaluiksi klassisten matemaattisten mallien rinnalle.
    Tavallisiin neuroverkkoihin verrattuna, fysiikkainformoitujen neuroverkkojen ennusteita pystyy rajoittamaan paremmin reunaehdoilla.
    Neuroverkkomallien laskennallinen tehokkuus mahdollistaa uusia nopeutta vaativia sovelluskohteita.  

    Havaintoihin ja malleihin liittyy usein epävarmuuksia.
    Kalman-suodin on yksi ratkaisu epävarmuuden arviointiin epävarmojen havaintojen ja mallien ongelmissa.
    Parven Kalman-suodin on tehokas Monte Carlo -tyyppinen approksimaatio, jota tarvitaan suurten datavektoreiden hallintaan.

    Tässä diplomityössä muodostetaan neuroverkko, joka tunnistaa advektioyhtälön liikekenttiä datasta.
    Neuroverkkomalli koulutetaan meterologisella pilvidatalla.
    Koulutusdatan vaikeuksita huolimatta, sillä koulutettu neuroverkko yleistyy odotetusti myös yksinkertaisempaan synteettiseen testidataan.
    Lisäksi neuroverkkomallin ennustuskyky on verrattavissa yleisesti käytettyyn optisen virtatuksen menetelmään.

    Neuroverkkoa testaan tilaoperaattorina parven Kalman-suotimessa.
    Testinä käytetään simuloitua tilannetta puuttuvien havaintojen täyttämiseksi.
    Neuroverkko toimii odotetusti myös Kalman-suodin mallin osana.
\end{abstractpage}
%% Force new page so that the Swedish abstract starts from a new page
\newpage

%% Swedish abstract. Delete it if you don't need it. 
%% 
\thesistitle{Fysikinformerade Neuronnät i Probabilistisk Rumstemporal Modellering}
\supervisor{TkD Harri Hakula}
\advisor{Prof.\ Marko Laine, FD Olle Räty} %
\degreeprogram{Mathematics and Operations Researach}
\major{Mathematics}
\department{Institutionen för matematik och systemanalys}%
%% Abstract keywords
\keywords{Fysikinformerade neuronnät, osäkerhetskvantifiering, Bayesisk filtrering}
%% Abstract text
\begin{abstractpage}[swedish]
Fysikinformerade neuronnät har visat sig vara värdefulla komplement till klassiska matematiska modeller.
Jämfört med vanliga fysikfria neuronnät kan randvillkor bättre upprätth\aa{}llas i de fysikinformerade varianterna.
Neuronnätens beräkningsmässiga effektivitet gör dem genomförbara för m\aa{}nga tillämpningar som inte är möjliga med klassiska metoder.

De flesta observationer i den verkliga världen och matematiska modeller är förem\aa{}l för osäkerheter.
Kalmanfiltret är en lösning för att kvantifiera osäkerhet i problem med osäkra observationer och modeller.
Ensemble-Kalmanfiltret ger en effektiv Monte Carlo-approximation när man hanterar stora datavektorer.

I denna avhandling konstrueras ett neuronnät för att härleda hastighetsfält inom advektion ekvationen.
Modellen tränas med meteorologiska molndata. Trots de inneboende sv\aa{}righeterna med träningsdatan visar sig neuronnätsmodellen kunna lära sig dynamiken i advektion ekvationen även med enklare syntetiska data.
Dessutom presenterar neuronnätet jämförbara resultat jämfört med en vanligt använd metod för skattning av optiskt flöde.

Nätverket används som en tillst\aa{}nds processoperator i ensemble Kalman filtret.
En simulerad inställning av partiella observationer med ett försök att producera kompletta tillst\aa{}ndsensembler testas.
Nätverket fungerar som förväntat inom Kalman filter modellen.
\end{abstractpage}

%% Force a new page after the preface
%%
\newpage


%% Table of contents. 
%%
\thesistableofcontents


%% \clearpage is similar to \newpage, but it also flushes the floats (figures
%% and tables).
%%
\cleardoublepage

%% Text body begins. Note that since the text body is mostly in Finnish the
%% majority of comments are also in Finnish after this point. There is no point
%% in explaining Finnish-language specific thesis conventions in English.
%% This text will be translated to English soon.
%%
\section{Introduction}
Neural networks have proved themselves an important tool in the analysis of high-dimensional and complex data being able to learn patterns and relationships.
In many physical modelling problems, existing mathematical models exist that characterise the evolution of a system but the models require parameters.
Direct observations of these model parameters often do not exist and must therefore be estimated from data.
For that challenge, neural networks could be constrained to learn the physical laws of a system to estimate parameters for the mathematical models.
These physics-informed neural networks are designed to leverage the prior information concerning a systems evolution to capture the essential patterns of even low amount of data.

However, in many problems, direct measurements from the quantity of interest are not available.
But even then, there may be partial measurements or measurments from one or more related quantities that provide useful information.
In these situtations, Bayesian filtering techniques such as the well-known Kalman filter provide a powerful method to probabilistically estimate the dynamic state of a system in terms of the actual quantity of interest.  

In this thesis, a neural network is constructed to predict advection fields from two-dimensional data, and it is used as a state model operator in an ensemble based Kalman filter.
The work is done as a part of a research project in the Meteorological Research Applications group at the Finnish Meteorological Institute.
Consequently, the model is trained and tested on meteorological satellite data of physical properties of clouds but the general model is designed to be independent of the used data.

In meteorological short-term forecasting, or nowcasting, the so-called physics-free models have been extensively researched in e.g.\ percipitation forecasts.~\cite{rainnet, metnet}
While these methods have shown some efficiency in forecasting applications, the neural networks often lack of interpretability and require large amounts of training data.
As actual physical numerical weather prediction models lack the ability to generate instanteneous high-resolution predictions, the currently used methods are crude mathematical models such as the advection based model studied in this thesis.
As a combination of the physics-free neural networks and the current advection based models, informing the neural network with the advection equation is explored. 
Similar work on physics-informed advection based neural networks have previously been done in~\cite{deepide, debezenac} with sea surface temperature data but in this work the objective is to generate smooth advection fields to better meet the requirements of small-scale cloud data. 
As already tested in~\cite{deepide}, the neural network model is integrated with the Kalman filter model but instead of the goal of probabilistic forecasting, the emphasis is on filling incomplete observations.

The main contributions of this thesis is the repurposing of an existing image recognition neural network architecture to a physics-informed neural network learning to produce smooth advection fields, and the integration of it to a high-dimensional ensemble based Kalman filtering scheme.
The scope of this work is primarily on the exploration and quick validation of new and alternative methods.
Consequently, the neural network development is mostly motivated by general data independent physical interpretability with less focus on performance based tuning.
Moreover, the Kalman filter is set up in an artificial setting in order to provide visually appropriate validation.  
Nevertheless, the neural network model is benchmarked against a well-used classical method used in advective meteorological nowcasting using a test partition of the used training dataset.

The thesis unfolds as follows.
First, some background material on the theory behind the used methods is presented.
Due to the wide range of the used methods and the theory behind them, naturally a complete primer is outside of the scope of this work.
As a compromise based on the spirit of this work, the background section is divided to three parts consisting of characteristics of spatio-temporal data, neural network basics, and Bayesian filtering theory. 
Even though the model was constructed hand in hand with the used dataset, due to it having been designed for general advective spatio-temporal data the proposed model is presented next independently. 
Finally, numerical experiments and verification of the trained model is presented with both the cloud data and synthetic test images separately for the neural network model and the Kalman filter application.
In the very last section, an overall summary of the work is discussed along with possible directions for further research from the basis of this work.


\clearpage
\section{Theoretical Background}
\input{background/background.tex}

\clearpage
\section{Model}
In this section, the hybrid model incorporating a physics-informed neural network with the Bayesian state-space modelling methods is proposed.
The objective is to create a framework where possibly multiple partial observations could be used to generate probabilistic estimates of some quantity of interest that evolves according to the advection equation.
The developed model is compromised of two main components: 
\begin{itemize}
    \item
        A deterministic neural network-based model that learns the patterns and dynamics of the advection equation from spatio-temporal data.
    \item
        An ensemble Kalman filter that leverages the neural network as a state evolution model to provide probabilistic state estimates.
\end{itemize}
These components are next presented separately, starting with the neural network model.

\subsection{Deterministic Neural Network Model\label{subsec:nnmodel}}
Departing from the previous research of the problem, a distinctive aspect of the model of this work is the generation of smooth advection fields as opposed to pixel-wise advection vector estimation.
The chosen approach is better suited especially for the high resolution meteorological applications in mind as small-scale differences in the data are mostly random given the available measurements, and features in the data mostly move to the same direction.

The data fed into the neural network consists of some smooth objects moving across two spatial dimensions over time.
Using several consecutive spatial snapshots of this data, the network predicts an advection field that represents the advection process across these time steps.
Even in the presence of noisy data, the objective is to produce smooth advectino fields.
The problem is depicted in figure~\ref{fig:advectionfigure}.
While the network is trained and evaluated mainly on satellite images of COT data, it is intended to be applicable to all similar processes.

Advection field estimation has also been studied in the context of optical flow problem.
Classical optical flow techniques are also often used in meteorology in nowcasting problems.
The neural network approach has also been studied but mostly focusing on different objectives.
They are often trained with ground truth flows and with datasets where motion only applies to small objects in images most notably with the KITTI~\cite{kitti}, and Sintel~\cite{sintel} datasets.
Contrary to this, the model used for this thesis has to be trained in an unsupervised fashion as the ground truth advection fields are not available, and the data for which the model is intended to be used with is very different from the usual ones used in the neural network optical flow context.

In order to tackle the needs of the much lower-dimensional data similar to figure~\ref{fig:advectionfigure}, a different kind of approach to the neural network structure is needed.
Basis vectors are utilised here akin to the method used in~\cite{deepide}, but a simple second degree polynomial is opted as opposed to a high-dimensional radial basis for the advection field.
The choice of outputting a low-dimensional polynomial basis from the neural network guarantees smooth advection fields. 

\subsubsection{Model Architecture}
The model architecture is similar to the well-known Residual Network (ResNet) that is deep yet computationally efficient.
The model compromises of two separate ResNet models, each trained to predict one component of the advection field: one for the vertical component, and the other for the horizontal component.

\subsubsection{ResNet Architecture}
Residual Networks are characterised by the use of skip connections, which allow the gradient to be directly backpropagated to earlier layers.
The residual term comes from these networks learning the residual mappings with reference to the layer inputs, alleviating the common vanishing gradient problem in deep neural networks.

A key component of the ResNet architecture is the residual block.
Residual blocks have the general form
\begin{equation}\label{eq:residualblock}
    F(x) + x,
\end{equation}
where $F(x)$ is the residual mapping to be learned representing layers of the neural network. $x$ is the original input to the block which represents the skip connection.
The important aspect here is the behaviour of the residual blocks in backpropagation.
Deep neural networks suffer from a problem called diminishing gradients when the chain rule results in a large number of multiplication of small numbers and thus producing small gradients.
However, differentiating the residual block gives
\begin{equation}\label{eq:residualblockderivative}
    \frac{d}{dx} F(x) + x = \frac{d}{dx} F(x) + 1.
\end{equation}
The $+1$ term ensures that even when propagating gradients through a deep network, they are not vanishing.

\subsection{Integrating the Neural Network into ETKF\label{subsec:etkfmodel}}
The neural network is trained to produce advection fields from spatio-temporal advective data but is only able to work with complete data and can only produce deterministic forecasts.
As real-world measurements of any process are rarely complete or without uncertainties, the neural network is used in the Bayesian filtering context to treat spatio-temporal data in the state-space model.

Due to the high dimensionality of image data, the ensemble transform Kalman filter of Algorithm~\ref{alg:etkf} is used to generate optimal state estimates from observations.
The concept of using the neural network model as a state evolution operator in an EnKF takes influence from~\cite{deepide} but here a slightly different approach is taken to accomodate for even larger state vectors and more realistic uncertainty estimates.

\subsubsection{State-Space Model Operators}
For analysis of the state-space model~\eqref{eq:statespacemodel}, observation and state evolution operators $H$ and $M$ are needed.
While the nonstationary inverse problem related to generating state estimates of measurements presents an important challenge in many problems, the main contribution of this work is in the incorporation of the neural network in the state evolution operator.
Therefore, the state and observation spaces are assumed to be of same dimension, and the observation operator is is simplified to be a linear identity mapping modified to filter out missing observations.

As the neural network models data through the advection equation by just generating advection fields, the state evolution operator is the linear warping function that transforms a state vector to the next time step
\begin{equation}
    M_{k+1}x_k = 
    \begin{cases}
        x_{k (i-\mathbf{F}_{u_{ij}}) ( j - \mathbf{F}_{vij})} & \text{if } \lvert i-\mathbf{F}_{u_{ij}} \rvert \leq n_x, \text{ and }  \lvert j-\mathbf{F}_{v_{ij}} \rvert \leq n_y
    \end{cases}
\end{equation}
Figure~\ref{fig:warping} depicts how the warping operator works.

\begin{figure}[H]
    \centering
    \includegraphics[width=5cm]{example-image-a}
    \caption{\label{fig:warping}The warping operator takes the pixel values and moves them to new pixels corresponding to an advection field.}
\end{figure}

\subsubsection{ETKF tuning}



\clearpage
\section{Numerical Experiments}
In this section the described model is trained and tested on meteorological cloud optical thickness (COT) data.
This section is divided to three parts.
Firstly, the cloud optical thickness data used for training is described.
Then the deterministic neural network model trained with the COT data and its performance evaluated.
And finally, the neural network is used as a parameter for the ETKF algorithm with the data.

\subsection{Cloud Optical Thicnkess Data}
While the described model is designed to be data-independent, in the development of it, Cloud Optical Thickness (COT) data retreived from satellites was used.
COT observations provide a consistent data source suitable for the model while really testing the capabilities of the model due to its qualities.
The choice of using COT data is motivated by two key factors; (1) it is easiliy available with regular capturing suitable for short-term forecasting applications, (2) the nature of the COT data makes it more or less suitable for the studied advection modelling.

Optical thickness is a quantity that is defined by the ratio of light passing through a material in logarithm, given by
\begin{equation*}
    \tau = \ln \left(\frac{I}{I_0}\right),
\end{equation*}
where $I$ is the transmitted light through the material and $I_0$ is the received light.
And in the cloud context, COT is a measure of a cloud's ability to absorb and scatter light, and it is an important parameter in satellite observations of clouds.

The data used for testing this model is computed using the EUMETSAT Nowcasting and Very Short Range Forecasting Satellite Application Facility (NWC SAF) cloud microphysics algorithms at the Finnish Meteorological Institute.
The algorithm computes the COT values using optical measurements from satellites on different wavelengths and a number of other parameters including atmosphere quantities computed using numerical weather prediction models.
For deeper understanding of the algorithmic principles behind the computation of the COT data, readers are referred to~\cite{nwcsaf}.

The data is captured with geostationary satellites that orbit 36000 km above the equator at a point where they remain at a constant position relative to the Earth.
The stationarity enables continuous observations of the same geographical region, which is indispensable for retrieving coherent spatio-temporal data.
However, due to the geostationary orbit above the equator, the spatial resolution of the data varies due to curvature of the earth and angle to the satellites.
Right under the satellite the spatial resolution is about \SI{3}{\kilo\metre}.
This dataset has a temporal resolution of 15 minutes, enabling accurate tracking of the evolution of clouds.

\subsubsection{COT Data Preprocessing}\label{subsec:cotpreprocessing}
In order to get usable COT data for the neural network, extensive preprocessing was required.
Advection modelling within the presented neural network necessitates complete spatio-temporal data in Euclidean coordinates.
The preprocessing of the incomplete raw COT data in a satellite projection is therefore a vital step before training the network.
As the network is relatively simple and should be trainable with a small dataset by its design, but preprocessing is computationally demanding, the data was spatially constrained to an area around southern Finland, and temporally to the months from April to October of 2021.

The spatial resolution of the data was artificially increased to have a pixel cover of approximately \SI{2}{\kilo\metre\squared}, thereby allowing clouds to occupy a larger image area and increasing the observable velocities for the network.
To accomodate for the neural network's square image inputs, the data was segmented into four $128 \times 128$ image regions.
Figure~\ref{fig:trainingdata_area} illustrates the image area and the four segments used for training the network.
\begin{figure}[h]
    \centering
    \inputpypgf{experiments/fig/traindata.pypgf}
    \caption{\label{fig:trainingdata_area}The area selected for neural network training over southern Finland and the Gulf of Finland.
    The entire image area is $512 \times 512$ pixels from which four $128 \times 128$ subimages were taken.
    One pixel covers roughly \SI{2}{\kilo\metre} in longitude and \SI{1}{\kilo\metre} in latitude.
    The missing measurements represented as black pixels present a major complication in usage of satellite data.}
\end{figure}

Satellite data often contain a considerable number of missing values as visualised in Figure~\ref{fig:trainingdata_area}.
As the missing values are mostly random with more occurence during some conditions including lower sun angles, realistic use of such data in forecasting use is thus problematic if needed at all times.
For training the neural network, some interpolation of the missing pixel values was mandatory to obtain a dataset of usable size.
Given the computational cost of more advanced interpolation, a simple spatial nearest valid value approach was employed.
While this crude method results in interpolation artefacts, for the purposes of the model training this is not necessarily a problem as the network is robust against poor and noisy data by its design.
As the dataset's quality is heavily influenced by the sun's elevation angle, interpolation was restricted to daytime images when the elevation angle exceeded \ang{15}.
As a last filter for the data, images with minimal COT values, indicating an absence of clouds, were subsequently filtered out from the dataset.
In this stage the data values were not yet normalised in any way to allow for quick experimentation within the neural network itself.
The finalised dataset was temporally segmented into distinct training, validation, and test sets laying the groundwork for subsequent neural network training.

A major difficulty posed by the COT data with regards to the advection modelling is the varying nature of its behaviour.
Large chunks of the finalised data consists of data that hardly evolves according to the simple advection equation.
While this could in principle have been bypassed by selecting only sequences of the data that evolve according to the advection model, it would have been difficult to achieve in practise and computationally restrictive.
However, the mixed quality of the training data tests the ability of the neural network to actually capture and learn the relevant information.

\subsection{Neural Network Model with COT Data}
The neural network presented in~\ref{subsec:nnmodel} was trained and tested on COT data.
\subsubsection{COT Data Preprocessing}
Advection modelling within the proposed neural network necessitates complete spatio-temporal data in Euclidean coordinates.
The preprocessing of the raw COT data is therefore a vital step before training the network.
As the network is relatively simple but preprocessing is computationally demanding, the data was spatially constrained to an area around southern Finland, and temporally to the months from April to October of 2021.

The spatial resolution of the data was artificially increased to have a pixel cover of approximately \SI{2}{\kilo\metre\squared}, thereby allowing clouds to occupy a larger image area and increasing the observable velocities.
To accomodate for the neural network's square image inputs, the data was segmented into four $128 \times 128$ image regions.
Figure~\ref{fig:trainingdata_area} illustrates the image area used for training the network.
\begin{figure}[h]
    \centering
    \inputpypgf{experiments/fig/traindata.pypgf}
    \caption{\label{fig:trainingdata_area}The area selected for neural network training over southern Finland and the Gulf of Finland.
    The entire image area is $512 \times 512$ pixels from which four $128 \times 128$ subimages were taken.
    One pixel covers roughly \SI{2}{\kilo\metre} in longitude and \SI{1}{\kilo\metre} in latitude.}
\end{figure}

Satellite data often contain missing values and thus interpolation was needed.
Given the computational cost of interpolation, a simple nearest valid value approach was employed.
As the dataset's quality is heavily influenced by the sun's elevation angle, interpolation was restricted to images when the elevation angle exceeded \ang{15}.

The processed data were then normalised to standard scores, narrowing the range of values to better align with preferences of convolutional neural networks.
Images with minimal COT values, indicating an absence of clouds, were subsequently filtered out from the dataset.

The finalised dataset was temporally segmented into distinct training, validation, and test sets laying the groundwork for subsequent neural network training.

\subsubsection{Training the Neural Network With COT Data}
The training process was tailored to the characteristics of the COT data while keeping in mind the underlying physical interpretation of advection fields.
The following choices were guided by performance of the model against the test dataset.

The preciously described architecture was chosen through extensive experimentations guided by performance, interprability, and model simplicity.
Within the ResNet model increasing the network depth provided little benefits and thus the network was kept quite simple.
The advection fields were produced using varioud basis types including Fourier, wavelet, and polynomial basis with different numbers of basis vectors.
The described second degree polynomials provided the best results in terms of advection field properties.
Furthermore, the image coordinates were placed within the positive real axes from 0 to 1 for evaluation of the polynomial functions.

Various loss functions were examined, including the MSE loss along with physics-based parameters accounting to the advection field properties such as gradient, curl, and divergence inspired by~\cite{debezenac}.
Also structural similarity (SSIM), multi-scale SSIM (MS-SSIM), logarithic cosh loss functions were tested.
Due to the small effect of the loss function in the results and no required parameters, the MSE loss was used for training the final model.

Three latest images proved to be the optimal amount of information for the network to produce advection fields.
The average loss was then computed over four future time steps.
Given the noisy nature of the data, a highest possible mini-batch size limited by GPU memory constraints was chosen so that bad training data examples wouldn't affect the training process.

Taking advantage of the advection fields' equivariance to rotations, the data was augumented by applying random rotations to the mini-batches.
With this, overfitting to the training data was avoided while permitting a smaller training dataset and shorter training times.



\subsection{Bayesian Filtering Results}


\clearpage
\section{Concluding Remarks}
In this work, an existing neural network architecture was modified for use in data independent physics-based advective modelling of spatiotemporal data and it was used as a state operator in a Kalman filter model.
The obtained neural network trained with a small dataset of satellite measurements of clouds was shown to generalise well for data unrelated to the training data while exhibiting comparable perormance to classic optical flow methods.
With using the Kalman filter with simple cloud data, the neural network was able to model the state evolution and thus fill missing observations.
In these experiments the concept of using neural networks in physical modelling was validated along with using it as a part of a more complex system.
While the data used in numerical experiments was limited to the cloud optical thickness observations, similar modelling could have applications across multiple domains.
The advective modelling was chosen for the generality and easy understanding but similar neural network frameworks could easily be applied also in modelling other phenomena where existing methods are not as well established.

While considerable work remains to follow up with what has been done for this thesis both with development of the neural network and tuning the Kalman filter is needed for better validation of the model, the results were already encouraging.
As the neural network was mostly a reused image recognition architecture with limited adjustments, a valid case could be made that there are considerable low-hanging fruits to improve performance.
For this work, the physics-based interpretability of the neural network model was a primary goal so engineering the network further with performance in mind and with more sophisticated training with extended and better appropriate data could easily provide superior results.
In application specific tasks, neural networks predicting outputs directly without physics-based modelling have gained popularity but lack in interpretability and in problems where mathematical models exist, the methods studied in this work could provide value in constraining the predictions.

In many ways the application of Kalman filtering in this work was limited and could not show its full potential.
While the examples studied here demonstrated its ability fill missing values in partial observations, more sophisticated treatment of nonstaionary inverse problems related to the observation operator in the state-space model could provide value in many complex problems.
For instance, data fusion of multiple data sources along with the methods studied here could contribute to meaningful improvements in accuracy of many problems.

Even with the wide use and theoretical validity of the ensemble based Kalman filters, major drawbacks remain to be handled.
The accurate representation of the uncertainty of state estimates requires handcrafted engineering of the ensembles with multiple unknown parameters.
But as larger state and observation spaces render the classical Kalman filter unfeasible, the ensemble methods offer a valid workaround.
Another alternative for applications similar to the one presented here could be appropriate dimension reduction techniques in order to continue using the classical Kalman filter for higher dimensional problems as well.
In addition, the filtering problems could be reformulated to better represent the actual task. 
As the neural network state operator models only advection at a few control points and big uncertainties are related to the advection vectors, a potential reformulation would be to incorporate the advection parameters to state vectors.

For meteorological applications, prediction problems related to the used cloud optical thickness observations could also be valuable even though less focus was given to that.
With reliable COT forecasts, the ever increasing solar energy production could be estimated that could be valuable information even for short time horizons in the electricity markets given the variable power of renewable energy sources.
For the forecasting problems, the presented model framework could be already used for complete and reliable estimates of the initial state from which with some adjustments probabilstic forecasts could be produced.

In conclusion, the findings of this work could direct further work toward multiple directions.
In the realm of the study of novel methods, similar techniques could be tested in similar modelling problems to the advection model.
On the other hand, this, and similar methods could benefit from futher verification and optimisation to yield better results.
As for the Bayesian filtering problems, the completing and smoothing of historical advective time series or the incorporation of more sophisticated observation models would be a direct continuation of this work.
Finally, the advective prediction problems related to a specific data and applications could be studied with the model of this work providing complete initial states and a predictor model.

\clearpage
\thesisappendix
\section{Proofs}\label{app:proofs}
\begin{proof}[Proof of Theorem~\ref{thm:updateformulas} following~\cite{kaipiosomersalo}]
    For the prediction step~\eqref{eq:predictionstep}, using the Markov property assumption~\eqref{eq:markov1} that $x_{k+1}$ is dependent only from $x_{k}$,
    \begin{align*}
        \pi(x_{k+1}, x_{k}, y_{1:k}) &= \pi(x_{k+1} \, \vert \, x_{k}, y_{1:k}) \pi(x_k, y_{1:k})\\
                                     &= \pi(x_{k+1} \, \vert \, x_{k}) \pi(x_k \, \vert \,  y_{1:k}) \pi(y_{1:k}).
    \end{align*}
    Substituting this to
    \begin{equation*}
        \pi(x_{k+1}, y_{1:k}) = \int \pi(x_{k+1}, x_{k}, y_{1:k}) \, d x_k,
    \end{equation*}
    and solving the conditional distribution $\pi(x_{k+1} \, \vert \, y_{1:k})$ gives the prediction step equation~\eqref{eq:predictionstep},
    \begin{equation*}
        \pi(x_{k+1} \, \vert \, y_{1:k}) = \frac{\pi(x_{k+1}, y_{1:k})}{\pi(y_{1:k})} = \int \pi (x_{k+1} \, \vert \, x_k) \pi (x_k \, \vert \, y_{1:k}) \, d x_k.
    \end{equation*}

    For the observation update step~\eqref{eq:updatestep},
    \begin{align*}
        \pi(x_{k+1} \, \vert \, y_{1:k+1}) &= \frac{\pi(x_{k+1}, y_{1:k+1})}{\pi(y_{1:k+1})}\\
                                           &= \frac{\pi(y_{k+1}, x_{k+1}, y_{1:k})}{\pi(y_{1:k+1})}\\
                                           &= \frac{\pi(y_{k+1} \, \vert \, x_{k+1}, y_{k}) \pi(x_{k+1}, y_{1:k})}{\pi(y_{1:k+1})}\\
                                           &= \frac{\pi(y_{k+1} \, \vert \, x_{k+1}) \pi(x_{k+1}, y_{1:k})}{\pi(y_{1:k+1})} \tag{\textasteriskcentered} \label{eq:updateformulas_star}\\
                                           &= \frac{\pi(y_{k+1} \, \vert \, x_{k+1}) \pi(x_{k+1} \, \vert \,  y_{1:k}) \pi(y_{1:k})}{\pi(y_{1:k}, y_{k+1})}\\
                                           &= \frac{\pi(y_{k+1} \, \vert \, x_{k+1}) \pi(x_{k+1} \, \vert \,  y_{1:k}) \pi(y_{1:k})}{\pi(y_{k+1} \, y_{1:k}) \pi(y_{1:k})}\\
                                           &= \frac{\pi(y_{k+1} \, \vert \, x_{k+1}) \pi(x_{k+1} \, \vert \,  y_{1:k})}{\pi (y_{k+1} \, \vert \, y_{1:k})}.
    \end{align*}
    Here~\eqref{eq:updateformulas_star} is justified by the conditional independence of $y_{k+1}$ from previous observations if $x_{k+1}$ is known.
\end{proof}

%\begin{proof}[Proof of Theorem~\ref{thm:kalmanfilterupdate} following~\cite{sarkka}]
    The predicted state mean results directly from the assumption of the model operator being an unbiased estimator.
    The prior distribution covariance requires studying a joint distribution.
    As the state and observation vectors are assumed to be Gaussian variables, so is the joint distribution of a pair of them.
    The joint distribution of $\pi(x_k \, \vert \, y_{1:k})$ and $\pi(x_{k+1} \, \vert \, y_{1:k})$ is
    \begin{align*}
        \pi\left( \begin{matrix}  x_k \, \vert \, y_{1:k}\\ x_{k+1} \, \vert \,  y_{1:k} \end{matrix} \right) \sim \mathcal{N}\left( \begin{bmatrix} x_{k\,\vert\,k} \\ M_{k+1} x_{k\,\vert\,k} \end{bmatrix},
          \begin{bmatrix}
              \Sigma_{11} & \Sigma_{12} \\
              \Sigma_{21} & \Sigma_{22}
          \end{bmatrix} \right).
    \end{align*}
    In the decomposed covariance matrix $\Sigma_{11} = P_{k \, \vert \, k}$, the covariance of $\pi(x_k \, \vert \, y_{1:k})$.
    $\Sigma_{12}$ is the covariance of $\pi(x_k \, \vert \, y_{1:k})$ and $\pi(x_{k+1} \, \vert \, y_{1:k})$.
    $x_{k+1} \, \vert \, y_{1:k} = M_{k+1} x_k \, \vert \, y_{1:k} + w_{k+1}$ by the Kalman filter assumptions where $w_{k+1} \sim \mathcal{N}(0, Q_{k+1})$.
    $\Sigma_{12}$ can be therefore evaluated using the affine equivariance of covariance and the linearity of expectation by
    \begin{align*}
        \Sigma_{12} &= \mathbb{E}[(x_k \, \vert \, y_{1:k} - \mathbb{E}[x_k \, \vert \, y_{1:k}])(x_{k+1} \, \vert \, y_{1:k} - \mathbb{E}[x_{k+1} \, \vert \, y_{1:k}])^{\top}]\\
                    &= \mathbb{E}[(x_k \, \vert \, y_{1:k} - \mathbb{E}[x_k \, \vert \, y_{1:k}])(M_{k+1}(x_k \, \vert \, y_{1:k} - \mathbb{E}[x_k \, \vert \, y_{1:k}]))^{\top}]\\
                    &= \mathbb{E}[(x_k \, \vert \, y_{1:k} - \mathbb{E}[x_k \, \vert \, y_{1:k}])(x_k \, \vert \, y_{1:k} - \mathbb{E}[x_k \, \vert \, y_{1:k}])^{\top}]M_{k+1}^{\top}\\
                    &=  P_{k \, \vert \, k}M_{k+1}^{\top}.
    \end{align*}
    $\Sigma_{21}$ is then just the transpose of $\Sigma_{12}$ by the symmetricity of the covariance matrix, $\Sigma_{21}= M_{k+1}P_{k \, \vert \, k}^{\top}$.
    $\Sigma_{22}$ is the covariance of $\pi(x_{k+1} \, \vert \, y_{1:k})$ and by using again the affine equivariance of covariance it can be written as
    \begin{align*}
        \Sigma_{22} &= \operatorname{Cov}(x_{k+1} \, \vert \, y_{1:k}, x_{k+1} \, \vert \, y_{1:k})\\
                    &= \operatorname{Cov}(M_{k+1} x_k \, \vert \, y_{1:k} + w_{k+1}, M_{k+1} x_k \, \vert \, y_{1:k} + w_{k+1})\\
                    &= M_{k+1} \operatorname{Cov}(x_k \, \vert \, y_{1:k}, x_k \, \vert \, y_{1:k}) M_{k+1}^{\top} + \operatorname{Cov}(w_{k+1})\\
                    &= M_{k+1} P_{k \, \vert \, k} M_{k+1}^{\top} + Q_{k+1} = P_{k+1 \, \vert \, k}.
    \end{align*}
    The distribution for $\pi(x_{k+1} \, \vert \, y_{1:k})$ can be given with the previous as
    \begin{equation*}
        \pi(x_{k+1} \, \vert \, y_{1:k}) \sim \mathcal{N}(M_{k+1}x_{k\,\vert\,k}, M_{k+1}P_{k\,\vert\,k}M_{k+1}^{\top} + Q_{k+1},
    \end{equation*}
    which concludes the first part.

    The observation update step is derived similarly through joint distributions.
    The joint distribution of the predicted state $\pi(x_{k+1} \, \vert \, y_{1:k})$ and the predicted observation $\pi(y_{k+1} \, \vert \, y_{1:k})$ is
    \begin{equation*}
        \pi \left(
            \begin{matrix}
                x_{k+1} \, \vert \, y_{1:k}\\
                y_{k+1} \, \vert \, y_{1:k}
            \end{matrix}
            \right) \sim \mathcal{N} \left(
                \begin{bmatrix}
                    x_{k+1 \, \vert \, k}\\
                    H_{k+1} x_{k+1 \, \vert \, k}
                \end{bmatrix},
                \begin{bmatrix}
                    \Sigma'_{11} & \Sigma'_{12} \\
                    \Sigma'_{21} & \Sigma'_{22}
                \end{bmatrix} \right).
    \end{equation*}
    With similar reasoning as above the covariance matrix consists of,
    \begin{align*}
        \Sigma'_{11} &= P_{k+1 \, \vert \, k}\\
        \Sigma'_{12} &= P_{k+1 \, \vert \, k} H_{k+1}^{\top}\\
        \Sigma'_{21} &= H_{k+1} P_{k+1 \, \vert \, k}^{\top}= H_{k+1} P_{k+1 \, \vert \, k}\\
        \Sigma'_{22} &= H_{k+1} P_{k+1 \, \vert \, k} H_{k+1}^{\top} + R_{k+1}.
    \end{align*}
    The normal conditional distribution $\pi(x_{k+1} \, \vert \, y_{k+1}) \sim \mathcal{N}(x_{k+1\,\vert\,k+1},  P_{k+1 \, \vert \, k+1})$ can be constructed from the above normal joint distribution (proof omitted).
    The expectation is given by
    \begin{align*}
        x_{k+1\,\vert\,k+1} &=\mathbb{E}[x_{k+1} \, \vert \,  y_{1:k}] + \Sigma'_{12} {\Sigma'_{22}}^{-1}(y_{k+1} - \mathbb{E}[ y_{k+1} \, \vert \, y_{1:k}])\\
                                                &= x_{k+1 \, \vert \, k} + P_{k+1 \, \vert \, k} H_{k+1}^{\top} (H_{k+1} P_{k+1 \, \vert \, k} H_{k+1}^{\top} + R_{k+1})^{-1} (y_{k+1} - H_{k+1}x_{k+1 \, \vert \, k})\\
                                                &= x_{k+1 \, \vert \, k} + K_{k+1}v_{k+1}.
    \end{align*}
    And the covariance is given by
    \begin{align*}
    P_{k+1 \, \vert \, k+1} &=\Sigma'_{12} - \Sigma'_{12} {\Sigma'_{22}}^{-1} \Sigma'_{21}\\
        &= P_{k+1 \, \vert \, k} - P_{k+1 \, \vert \, k} H_{k+1}^{\top} (H_{k+1} P_{k+1 \, \vert \, k} H_{k+1}^{\top} + R_{k+1})^{-1} H_{k+1}P_{k+1 \, \vert \, k}\\
        &= P_{k+1 \, \vert \, k} - K_{k+1} H_{k+1}P_{k+1 \, \vert \, k}\\
        &= (I -  K_{k+1} H_{k+1}) P_{k+1 \, \vert \, k}.
    \end{align*}
    This concludes the proof.
\end{proof}



%% Leave page number of the first page empty
%% 
%\thispagestyle{empty}

\clearpage
%% L\"ahdeluettelo

%\thesisbibliography
\bibliographystyle{plain}
\bibliography{references}


\end{document}
